% Generated draft: requires mathematical and source review.
\documentclass[11pt]{article}
\usepackage[margin=25mm]{geometry}
\usepackage{fontspec}
\setmainfont{Latin Modern Roman}
\usepackage{amsmath,amssymb,mathtools,unicode-math}
\setmathfont{Latin Modern Math}
\usepackage{enumitem,needspace,xurl,hyperref}
\setlength{\parindent}{0pt}
\setlength{\parskip}{5pt}
\setlength{\emergencystretch}{3em}
\newcommand{\N}{\mathbb{N}}
\newcommand{\Z}{\mathbb{Z}}
\newcommand{\Q}{\mathbb{Q}}
\newcommand{\R}{\mathbb{R}}
\newcommand{\C}{\mathbb{C}}
\newcommand{\F}{\mathbb{F}}
\newcommand{\A}{\mathbb{A}}
\newcommand{\PP}{\mathbb{P}}
\newcommand{\E}{\mathbb{E}}
\newcommand{\eps}{\varepsilon}
\newcommand{\dd}{\,\mathrm{d}}
\begin{document}
\section*{UnsolvedMath ID: 30002169}
\textbf{Source problem code:} catalogue entry
\subsection{Definitions and mathematical statement}
The supplied problem concerns rough Hamilton--Jacobi--Bellman equations driven by a rough path \(\mathbf z\), understood as a driving signal equipped with the iterated-integral information required by the intended rough-path theory. The source record does not specify the equation, the roughness regime, the spatial domain, the time interval, the coefficients, or the initial and boundary data. These must be fixed before a fully quantified mathematical problem can be stated.

For such a fixed equation and data, let \(z^{(n)}\) be smooth driving signals whose rough-path lifts \(\mathbf z^{(n)}\) converge to \(\mathbf z\) in the intended rough-path topology. Let \(u^{(n)}\) be viscosity solutions of the corresponding equations with smooth driving signals and the prescribed data. An approximation-defined solution is a limit
\[
\mathbf z^{(n)}\longrightarrow\mathbf z,
\qquad
u^{(n)}:=u^{(n)}\longrightarrow u,
\]
where the topology of solution convergence and the admissibility requirements on the approximations remain unspecified in the supplied record. No existence, uniqueness, or independence of the approximating sequence is asserted here.

\textbf{Problem.} Find an intrinsic characterization of these approximation-defined solutions: formulate a condition on \(u\), the rough driver \(\mathbf z\), the equation, and its data that is equivalent to being a solution in the intended approximation-based sense, without referring to a chosen sequence of smooth driving signals. In particular, can this be achieved by a viscosity-type definition using suitable test functions that touch \(u\) from above or below, together with appropriate subsolution and supersolution conditions?

The intended characterization should apply to every driver and candidate solution in the intended equation class. A touching-test formulation would need to specify the test-function class, its dependence on the rough driver, the meaning of touching, and the inequalities imposed at contact points. The source does not supply these specifications, nor does it state whether approximation-based solvability means convergence along some admissible sequence or a common limit along every admissible sequence. Consequently, this is a statement of the requested characterization problem, with unresolved parameters made explicit, rather than a fully specified PDE conjecture.
\subsection{Short English statement}
Can solutions of rough Hamilton--Jacobi--Bellman equations obtained by smoothing the driving signal and taking limits be recognized directly from the rough driver and the candidate solution, for example through suitable touching test functions? The equation class, convergence rules, and admissible tests require specification. The supplied record labels the problem open; that status has not been independently verified.
\subsection{Sources}
\begin{itemize}
\item Supplied local UnsolvedMath registry record, permanent ID \(30002169\), exact source problem code \texttt{catalogue entry}, in the collection \emph{Oberwolfach Reports Open Problems}. No individual proposer is identified; the proposed year is recorded as 2012.
\item Supplied report citation: \emph{Rough Paths and PDEs (2013)}, Oberwolfach Reports, DOI \texttt{10.4171/owr/2012/41}, pp.\ 2493--2540. \url{https://doi.org/10.4171/owr/2012/41}. This is also the supplied literature source labelled ``OWR report 41.'' The background describes Oberwolfach Reports as documenting mathematical workshops and their open-problem discussions.
\item Supplied literature source labelled ``Rough PDE viscosity theory'': \url{https://doi.org/10.1186/s41546-020-00049-8}. No author names or fuller bibliographic details are supplied for this reference.
\item The exact supplied literature-check date is \texttt{2026-08-21}. The supplied dated curation labels the problem open and assesses that pathwise viscosity and controlled-test-function theories handle important rough and stochastic PDE classes, while a general intrinsic touching-test characterization remains unresolved. The record expressly describes this as dated literature triage derived from the supplied Oberwolfach curation and calls for independent verification of time-sensitive status and cited claims. This assessment is supplied dated curation, not independently verified here.
\item The supplied statement-status label is \texttt{corrected\_verified}. Its accompanying verification note states that the official report explicitly asks for an intrinsic characterization of the pathwise viscosity solution. This is a supplied source-verification claim; the report has not been independently consulted here.
\end{itemize}
\end{document}
